\chapter{The exact obligation that remains}\label{ch:closeout}
\section{The target after every valid reduction}
The constructions in this reader return to one fixed arithmetic question. They have not replaced $\Delta$ by a more convenient sequence, and they have not changed nonvanishing into a statement about typical primes. Their target is still every actual prime coefficient.

\begin{proposition}[The remaining family]\label{prop:remainingfamily}
Lehmer's conjecture is equivalent to the assertion
\begin{equation}\label{eq:remainingfamily}
R(p)=\tau(p)\ne0
\qquad\text{for every prime }p\equiv4\pmod5.
\end{equation}
\end{proposition}
\begin{proof}
Theorem~\ref{thm:primereduction} reduces the conjecture to prime indices. The congruence \eqref{eq:primemod5} excludes a zero at primes congruent to $1,2,3$ modulo $5$, and the explicit value $\tau(5)=4830$ handles $5$. The primitive transform equals $\tau$ at a prime by \eqref{eq:rprime}. Thus exactly the displayed unexcluded family remains.
\end{proof}

The word \emph{remaining} describes the family after these universal reductions. Many of its individual coefficients have also been certified, both in the finite examples here and in the much larger historical computation cited in Chapter~\ref{ch:target}. Accepting the latter as an external input lets one restrict attention to primes at or above its stated endpoint. It does not remove the universal prime quantifier.

On each such prime shell, all packets have content one. The residual is therefore the full signed sum
\begin{equation}\label{eq:primepacketobligation}
\tau(p)=\sum_{v\in\mathcal V_p}W(v).
\end{equation}
The summands are nonzero integers, and for $p>20$ they are units at the index prime. Neither property forbids their sum from being zero. The unresolved statement concerns the exact balance of their signed masses.

\section{What the program has proved}
The side theorems are useful precisely because their quantifiers and dependencies are now explicit. Primitive extraction has a complete local description. It gives a nonzero residual at every square and introduces no new zeros anywhere. The polynomial construction supplies universal null corrections, with harmonicity and valence as the justification. Finite rational sign certificates prove their individual coefficients survive. The tree theorem identifies the first mixed-sign shell at every dyadic depth. Content conservation persists for every finite symmetric polynomial universal null, without a degree cutoff or a homogeneity assumption.

\Needspace{15\baselineskip}
\begin{center}\small
\begin{tabularx}{\linewidth}{>{\raggedright\arraybackslash}p{0.25\linewidth}X}
\toprule
Established statement&Exact implication for the target\\
\midrule
$R(s^2)\ne0$ for all $s$&Every even primitive local block survives, including a hypothetical zero prime parameter.\\
$R(n)=0\iff\tau(n)=0$&Primitive extraction encodes the same zero problem, with $R(p)=\tau(p)$.\\
Universal polynomial nulls&Permitted corrections preserve every coefficient and every exact-content mass.\\
Sharp dyadic boundary&Geometric separation is controlled at all depths, without comparing the global positive and negative totals.\\
Exact finite certificates&The stated shells, matrix ranks and polynomial residues are settled in their declared ranges.\\
\bottomrule
\end{tabularx}
\end{center}

These statements are written proofs and exact finite certificates. They are not a compiled proof-assistant solution of Lehmer. Their inclusion claims no historical priority. In particular, a translated classical identity and a project-specific consequence of that identity have different provenance even when both are correct.

The arithmetic survivor theorem also has an all-modulus quantifier, but it is a limitation theorem for a proposed method. For every fixed positive modulus $M$, a positive-density set of actual primes satisfies
\[
p\equiv-1\pmod{5M},
\qquad R(p^r)\equiv0\pmod M\quad(r\ge1).
\]
It applies to the actual joint finite image, with no independence assumption. This is stronger than observing survivors in a bounded search. It still says nothing about an integer zero at one fixed prime, because that prime may change with $M$.

\section{The missing arrows cannot be supplied by repetition}
Several attractive final steps have failed in precisely stated forms. The integral prime-index interpolation violates a necessary divided-difference condition. The aggregate auxiliary polynomial does not inherit its individual factors' root interval. The proposed lifting equalities at $2,3,5$ have exact prime counterexamples. Strict preservation of sectors with opposite masses cannot yield a common-sign sum. Interpolation on one shell does not give universal nullness, and top-degree harmonic projection does not preserve the interpolated values.

Other statements are unproved rather than refuted. The all-degree universal-null space has not been shown to fill every content-preserving transfer space at every shell. Its finite rank profiles are examples, not a uniform theorem. Even a proof of that fullness would retain the question of whether the conserved prime total is zero.

The local filtered Frobenius construction and the arithmetic sequence $f$ test different kinds of retained information. The former realizes compatible local zero-trace data. The latter satisfies a specified collection of arithmetic rules while having a prime zero. Neither is identified unconditionally with an actual zero coefficient of $\Delta$. Their role is to show why those weaker premises alone cannot force nonvanishing.

A new attempt should therefore state which missing arrow it supplies. Rechecking the recurrence, increasing a coefficient endpoint or adding finitely many fixed congruence tests does not create that arrow. This is a logical conclusion about the retained premises. It does not declare every possible use of recurrences, computation or congruences exhausted.

\section{A sufficient adaptive residue input}
One precise route would be a uniform construction of residue certificates. Suppose there are explicitly defined functions on prime indices, with integer values $m(p)\ge2$ and $b(p)$, for which one proves
\begin{equation}\label{eq:adaptiveresidueinput}
\tau(p)\equiv b(p)\pmod{m(p)},
\qquad b(p)\not\equiv0\pmod{m(p)}
\end{equation}
for every prime in \eqref{eq:remainingfamily}. These two assertions would settle that family and hence Lehmer. The modulus may depend on $p$, so the fixed-modulus survivor theorem does not refute this proposed form of input.

The issue is the proof of the second assertion for the actual coefficient, together with the identity linking the construction to $\tau$. Defining $b(p)$ to be the unknown residue and then assuming it is nonzero would only restate the target. A useful construction must give an independently justified nonzero residue, with all exceptional primes and all precision requirements covered.

The index-dependent criterion already available from Deligne's bound is
\[
\tau(p)=0\iff p^6\mid\tau(p),\qquad p\ge5.
\]
Combined with the divisor convolution, it becomes
\begin{equation}\label{eq:closeoutconvolution}
840S_4(p)\not\equiv p^4C(p)\pmod{p^6},
\qquad p>7,
\end{equation}
as the exact desired exclusion. This gives a correctly typed target for a convolution route. No uniform arithmetic property excluding that residue has been proved here. A stronger valuation bound may suffice, but its proof must permit the zero case as a possibility until the bound itself excludes it.

\section{A sufficient sign input}
A second precise route would construct permitted corrections with an independently proved sign. At each remaining prime, choose finitely many rational universal null polynomials $Q_{p,j}$ and rational coefficients $a_{p,j}$, and put
\[
F_p(v)=1+\sum_j a_{p,j}Q_{p,j}(v).
\]
The choice may depend on $p$. Universal nullness must hold after the choice, so its weighted sum is still the actual $\tau(p)$.

Suppose the construction proves a sign $\varepsilon_p\in\{-1,1\}$ for which
\begin{equation}\label{eq:uniformsigninput}
\varepsilon_p W(v)F_p(v)\ge0
\quad\text{for every }v\in\mathcal V_p,
\end{equation}
and proves that at least one of these finitely many quantities is strictly positive. Then summing \eqref{eq:uniformsigninput} gives $\varepsilon_p\tau(p)>0$. This would be a valid all-prime certificate theorem.

The strict clause cannot be discarded. If the total were zero, a correction making every term zero would satisfy the weak inequalities. Nor can the strict clause be obtained merely from full transfer rank. Equalizing corrected masses to $\tau(p)/N$ already uses the total whose nonvanishing is at issue.

The finite certificates at $7$ and $31$ show that this form of argument can work at concrete shells. A uniform construction needs a proof covering all remaining primes, with coefficients and inequalities derived from actual packet structure. The current family and its finite successes do not provide such a theorem.

\section{Where a global input would have to enter}
The actual form has more structure than an arbitrary multiplicative recurrence sequence. A new global theorem could impose a relation or constraint that is absent from the countermodel and that excludes a zero prime parameter. The constraint must be stated exactly and derived for $\Delta$. Merely naming modularity or compatibility does not prove a particular trace is nonzero.

Local compatibility also has a direction. A weakly admissible trace-zero object establishes that the selected local conditions permit zero. It does not establish that this object occurs as the restriction of the actual global representation. Proving such an identification at a prime, together with a proved comparison identifying its normalized Frobenius trace with $\tau(p)$, would give an actual zero and would disprove Lehmer. Proving a global obstruction to that identification could instead contribute to nonvanishing. Neither direction has been supplied by the local construction alone.

Similarly, showing that the synthetic series $\sum_{n\ge1}f(n)q^n$ is a holomorphic cusp form of level one and weight $12$ would be new information. Its normalization $f(1)=1$ and the one-dimensional cusp space would then identify it with $\Delta$. The present construction does not establish that premise. Calling the sequence nonmodular without an independent argument would also go beyond what was proved, since it could coincide with $\tau$ if the chosen prime were already an actual zero.

An effective continuation needs a concrete additional input, its exact hypotheses and a proof that it changes the zero test. That input may concern a prime-shell identity, an adaptive arithmetic restriction or a global relation. This reader establishes none of those prospective finishing theorems.

\section{The closeout statement}
The publication preserves both positive results and failed implications. A failed shortcut is identified by its exact domain and counterexample. An open span or global realization statement remains open. A finite run retains its finite range. A universal side theorem retains the proof that licenses its universal quantifier.

\begin{openproblem}[One actual coefficient at every remaining prime]
For every prime $p\equiv4\pmod5$, prove that the actual signed prime-shell total in \eqref{eq:primepacketobligation} is nonzero. Supply an independent nonzero residue, a legitimate sign certificate with a strict term, or another fully justified mechanism excluding exact cancellation. A verified actual prime zero would instead disprove the conjecture.
\end{openproblem}

The program has clarified the objects, their conserved quantities and several boundaries on their use. Its final mathematical status is that the prime total is still unresolved. The source and evidence make this unresolved status inspectable.
